\section{Introduction}
\label{sec:introduction}

An arithmetic partition function can resemble a dynamical zeta expression
without carrying the primitive objects that a dynamical determinant counts.
The distinction is especially sharp for the rooted Knauf binary recursion:
its limiting partition sum is an exact quotient of zeta functions, while its
finite source states are rooted words embedded by trailing zeroes.  The
quotient identity and its phase-transition setting belong to the original
number-theoretic spin-chain literature
\citep{knauf1993phases,knauf1998spin,knauf1999erratum}.

The issue studied here is not whether that arithmetic identity is valid.  It
is whether the very same finite object and label also determine the data
needed for a scalar primitive-orbit determinant: a state on which the binary
step acts autonomously, a clock that is independent of the root of a cyclic
word and additive under repetition, a scalar phase with the corresponding
power law, and an operator whose trace powers count those same returns.  A
partition trace supplies none of these identifications by notation alone.

We therefore evaluate the smallest object-preserving bridge.  We retain the
source matrices, the rooted label
$h(w)=\onevec^{\mathsf T}M_we_1$, the stable-state quotient under trailing
zeroes, ordinary cyclic rotation, ordinary word powers, and the literal
Liouville observable $\lambda(h(w))$, where
$\lambda(n)=(-1)^{\Omega(n)}$ and $\Omega(n)$ counts prime factors with
multiplicity.  Exact words of length at most three then settle the necessary
descent questions; \Cref{fig:rooted-state} previews the object-preserving
bridge and its first failure.  A separate trace-class
argument identifies the diagonal determinant that the stable-state inventory
actually owns.  Neither argument uses target zeros, fitted parameters, or a
prime-indexed external component.

Related Farey constructions make the typing issue concrete.  A
translation-invariant Farey chain can share free energy with a related
non-translation-invariant number-theoretic chain, and generalized Farey
models connect matrix recursions to intermittent transfer operators
\citep{klebanOzluk1999farey,fiala2003phase,fialaKleban2005generalized}.
Trace-energy chains exploit cyclic invariance, and prescribed-trace matrix
products have their own arithmetic counting theory
\citep{prellbergFialaKleban2006cluster,technau2023farey}.  These are genuine
models, but equality of free energy or use of the same generators does not
identify their finite objects, clocks, or return operators with the rooted
first-column-sum object tested here.

The same firewall applies to determinant language.  Orbit zeta and
transfer-operator determinants are genuine dynamical constructions when
their phase spaces, weights, and operator domains are specified
\citep{ruelle1976zeta,mayer1990gauss,mayer1991selberg}; their existence does
not identify an operator for the rooted label.  Conversely, the diagonal
operator constructed below is a valid Fredholm object, but its free marker
enumerates powers of inventory eigenvalues rather than binary returns.

The contribution is therefore a closure theorem, not a new zeta mechanism.
It has three parts.  First, the canonical append-one step does not descend to
the trailing-zero colimit.  Second, the rooted clock and literal scalar phase
fail cyclic and repetition descent by explicit witnesses.  Third, the exact
state-inventory determinant is derived and assigned to its legal owner.  The
source recursion and quotient, thermodynamic trace-chain repairs,
generalized transfer operators, and the general partition-trace principle
remain prior art.  The selection that led to this test is retrospective and
receives no novelty or priority credit.

\Cref{sec:prior-scope} records this source ownership and the claim boundary.
\Cref{sec:source-types} fixes the types and matrix convention.
\Cref{sec:non-descent} proves the finite non-descent theorems, and
\Cref{sec:inventory} derives the determinant owned by state inventory.
\Cref{sec:route} maps those results to the strict Route coordinates without
adding experiment claims.  Complete exact products, domain restrictions, and
literature firewalls appear in the appendices.
